\subsection{Boundary comparison for actual harmonic functions}\label{subsec:fc-line-harmonic-maximum}

\begin{theorem}[Weak maximum on the whole finite boundary]\label{thm:fc-harmonic-weak-maximum}
\lean{EnvelopeProperty.Berkovich.RelativeSpectrum.FCStandardAffinoid.harmonicExpression_exists_boundary_max, EnvelopeProperty.Berkovich.RelativeSpectrum.FCStandardAffinoid.harmonicExpression_exists_boundary_ge}
\leanok
\uses{def:fc-actual-harmonic, lem:fc-standard-inflation,
  lem:fc-standard-boundary, thm:fc-harmonic-local-maximum}
Every actual harmonic expression on a standard affinoid attains its maximum
at one of the actual points of the finite Gauss boundary. Consequently, for
every $x\in U$ there is $z\in\partial U$ with $h(x)\le h(z)$.
\end{theorem}

\begin{proof}
The actual carrier is nonempty and compact, and the expression is continuous
on it, so a maximum $M$ is attained. Suppose no boundary point attains $M$.
Every point of the maximum locus is then in the topological interior of $U$.
At such a point the global maximum inequality on $U$ is an actual local
maximum inequality in the analytic line. Coordinate positivity follows from
the positive hole radii. The all-point local maximum theorem therefore makes
the expression constant on an actual neighbourhood. Thus the maximum locus
is open in $U$. It is closed in $U$ by continuity and nonempty by attainment.
Since $U$ is connected, this locus is all of $U$. It contains the outer Gauss
boundary point, contradicting the supposition. A boundary maximum exists,
and comparison with it gives the asserted inequality for every point.
No distinctness or positive-length condition on boundary segments is used.
\end{proof}

\begin{corollary}[Zero boundary and Dirichlet uniqueness]\label{cor:fc-harmonic-boundary-uniqueness}
\lean{EnvelopeProperty.Berkovich.RelativeSpectrum.FCStandardAffinoid.harmonicExpression_eq_zero_of_boundary_zero, EnvelopeProperty.Berkovich.RelativeSpectrum.FCStandardAffinoid.harmonicExpression_eq_of_boundary_eq, EnvelopeProperty.Berkovich.RelativeSpectrum.FCStandardAffinoid.IsHarmonic.eq_of_boundary}
\leanok
\uses{thm:fc-harmonic-weak-maximum}
A harmonic function with zero values on the actual boundary is zero on $U$.
Two actual harmonic functions with the same boundary values are equal.
\end{corollary}

\begin{proof}
Apply the weak maximum principle to the function and its negative. Both are
at most zero on $U$, so the function is zero there. The difference of two
harmonic expressions has coefficients given by the differences of their
constants and finite coefficient vectors, hence is harmonic. Apply the zero
boundary result to their difference. This proves equality of the actual
functions and does not assert uniqueness of coefficient witnesses.
\end{proof}

\begin{theorem}[Strong maximum on an interior component]\label{thm:fc-harmonic-component-maximum}
\lean{EnvelopeProperty.Berkovich.RelativeSpectrum.FCStandardAffinoid.harmonicExpression_eq_on_interior_component}
\leanok
\uses{thm:fc-harmonic-local-maximum, lem:fc-standard-boundary,
  lem:fc-harmonic-retraction}
If a harmonic expression has its global maximum on $U$ at an interior point
$x$, it is constant on the connected component of $x$ in
$U\setminus\partial U$.
\end{theorem}

\begin{proof}
Restrict the expression to the open interior. Its maximum locus is relatively
closed by continuity. At every point of that locus the global maximum
inequality gives an actual local maximum in the analytic line; the all-point
local maximum theorem makes the locus relatively open as well. It is a
clopen subset of the interior containing $x$. Every connected component
meeting a clopen subset lies in that subset. Thus the whole interior
component of $x$ has the same value. This formulation permits boundary
points which are internal vertices of the skeleton.
\end{proof}
